\clearpage
\phantomsection
\chapter*{{\khmoul អត្ថបទសង្ខេប}}
\addcontentsline{toc}{chapter}{{\khmoul អត្ថបទសង្ខេប}}

\begin{khmer}

\hs ការត្រួតពិនិត្យហេដ្ឋារចនាសម្ព័ន្ធសំខាន់ៗ ដូចជា ស្ពាន អគារ និងផ្លូវរត់យន្តហោះ មានសារៈសំខាន់ណាស់សម្រាប់សុវត្ថិភាពសាធារណៈ។ ទោះជាយ៉ាងណា វិធីសាស្ត្រត្រួតពិនិត្យដោយដៃនាពេលបច្ចុប្បន្នមានចំណាយខ្ពស់ ពឹងផ្អែកលើការវាយតម្លៃផ្ទាល់ខ្លួន និងមានដែនកំណត់ក្នុងការគ្របដណ្តប់។ ទោះបីជាបច្ចេកវិទ្យា Deep Learning ទទួលបានជោគជ័យក្នុងការស្វែងរកស្នាមប្រេះលើរូបភាពថតជិតក៏ដោយ ការអនុវត្តលើរូបភាពអាកាសចរណ៍កម្រិតច្បាស់ខ្ពស់ដែលថតដោយដ្រូន (UAV) ប្រឈមនឹងបញ្ហាធំៗ៖ ស្នាមប្រេះដែលមានទទឹងត្រឹម ១-៣ ភិកសែល ត្រូវបង្ហាញលើរូបភាពទំហំជាង ១០០ មេហ្គាភិកសែល ដែលមានផ្ទៃតិចជាង ០.០១\% នៃរូបភាពសរុប។ ទំហំដ៏ធំនេះលើសពីសមត្ថភាពអង្គចងចាំរបស់ GPU ស្តង់ដារ និងធ្វើឱ្យគំរូទូទៅបាត់បង់ព័ត៌មានលម្អិត ឬបរិបទរួម។

\hs គម្រោងសិក្សាស្រាវជ្រាវរយៈពេល ៣ ខែនៅក្រុមហ៊ុន AI Farm Robotics Co., Ltd. បានបង្កើត និងផ្ទៀងផ្ទាត់ប្រព័ន្ធ Deep Learning សម្រាប់កំណត់ទីតាំងស្នាមប្រេះកម្រិតភិកសែលដោយស្វ័យប្រវត្តិលើរូបភាពអាកាសចរណ៍ទំហំធំ។ ការស្រាវជ្រាវបានប្រៀបធៀបស្ថាបត្យកម្មគំរូភាគរយចំនួនបី៖ U-Net (ប្រើ EfficientNet-B0 encoder និង scSE attention), SegFormer (ប្រើ MiT-B0 encoder) និង YOLO26s-seg។ U-Net និង SegFormer បានប្រើប្រាស់មុខងារបាត់បង់រួម (BCE + Dice + Focal Loss) ដើម្បីដោះស្រាយបញ្ហាភាពខុសគ្នាច្រើនរវាងផ្ទៃខាងក្រោយ និងស្នាមប្រេះ ខណៈដែល YOLO26s-seg ត្រូវ បានប្រើជាគំរូមូលដ្ឋានល្បឿនលឿន។ លើសពីនេះ ប្រព័ន្ធដំណើរការរូបភាពតាមបំណែក (Tile-based inference) ដែលមានការជាន់គ្នា ២៥\% ត្រូវបានរចនាឡើងដើម្បីអាចដំណើរការរូបភាពទំហំធំដោយមិនកម្រិតលើ GPU តែមួយ។

\hs តាមរយៈការបណ្តុះបណ្តាលលើរូបភាពស្នាមប្រេះសាធារណៈចំនួន ៨,៨១៨ រូបភាព (បែងចែក ៨០\% បណ្តុះបណ្តាល, ១០\% ផ្ទៀងផ្ទាត់, ១០\% សាកល្បង) U-Net ទទួលបានលទ្ធផលត្រឹមត្រូវបំផុតលើសំណុំទិន្នន័យសាកល្បង (IoU ៥៣.០៤\%, F1 ៦៧.៥៩\%)។ ផ្ទុយទៅវិញ YOLO26s-seg មានល្បឿនលឿនជាងគេ (១៨.០៧ ms) ប៉ុន្តែទទួលបានភាពត្រឹមត្រូវទាបជាង (IoU ២៨.១០\%)។ ក្នុងការសាកល្បងលើទិន្នន័យ UAV ថ្មី (CUBIT-InSeg) ដោយមិនបានបណ្តុះបណ្តាលមុន (Zero-shot) ប្រសិទ្ធភាពនៃគំរូទាំងអស់បានធ្លាក់ចុះ ដោយ YOLO26s-seg ទទួលបានលទ្ធផលល្អជាងគេបន្តិច (IoU ៤០.១២\%, F1 ៥៦.១៦\%)។ បច្ចេកទេសបែងចែករូបភាពតាមបំណែកបានជួយរក្សាភាពភ្ជាប់គ្នានៃស្នាមប្រេះឆ្លងកាត់ព្រំដែនបំណែករូបភាពបានយ៉ាងល្អ។ លទ្ធផលទាំងនេះបញ្ជាក់ថា U-Net ជាគំរូមានមូលដ្ឋានរឹងមាំសម្រាប់កំណត់ស្នាមប្រេះលម្អិត ហើយការស្រាវជ្រាវនេះបានផ្តល់គ្រឹះយ៉ាងសំខាន់សម្រាប់ការអភិវឌ្ឍបន្ថែមលើការត្រួតពិនិត្យហេដ្ឋារចនាសម្ព័ន្ធជាក់ស្តែង។ 

\vspace{0.5cm}
\noindent\textbf{ពាក្យគន្លឹះ:} ការបែងចែកស្នាមប្រេះ រូបភាព UAV រៀនសូត្រជ្រៅ U-Net SegFormer YOLO សន្និដ្ឋានផ្អែកលើក្បឿង ភាពមិនសមតុល្យថ្នាក់ ការត្រួតពិនិត្យហេដ្ឋារចនាសម្ព័ន្ធ

\end{khmer}